\begin{vita}

\dateitem{2013}{B.Tech Computer Science, VNRVJIET, Hyderabad, India}

\dateitem{2018}{M.S. Computer Science, University of Minnesota - Duluth, Minnesota, USA}



\begin{publist}

%% UPDATE FOR 2010:
%  Grad school only wants research publications, and it only wants those
%  research pubs that are actually published. Accepted or ``to appear''
%  publications don't count. If they look closely, they'll tell you to
%  remove any publications that aren't in print. Haivng said that, they
%  probably won't look that closely unless you put a really long list
%  here. You're tempting fate if you add instructional publications
%  though.

\researchpubs

\pubitem{M.~S.~Vallabhajosyula and R.~Ramnath
\newblock ``Towards Practical, Generalizable Machine-Learning Training Pipelines to Build Regression Models for Predicting Application Resource Needs on HPC Systems''.
\newblock {\em Practice and Experience in Advanced Research Computing 2022: Revolutionary: Computing, Connections, You}, pp. 1--5, 2022.}

\pubitem{S.~Vallabhajosyula and R.~Ramnath
\newblock ``Establishing a Generalizable Framework for Generating Cost-Aware Training Data and Building Unique Context-Aware Walltime Prediction Regression Models''.
\newblock {\em 2022 IEEE Intl Conf on Parallel \& Distributed Processing with Applications, Big Data \& Cloud Computing, Sustainable Computing \& Communications, Social Computing \& Networking \\(ISPA/BDCloud/SocialCom/SustainCom)}, pp. 497--506, 2022.}


\pubitem{M.~S.~Vallabhajosyula and R.~Ramnath
\newblock ``Towards Characterizing DNNs to Estimate Training Time using HARP (HPC Application Resource (runtime) Predictor)''.
\newblock {\em Practice and Experience in Advanced Research Computing 2023: Computing for the Common Good}, pp. 483--485, 2023.}

\pubitem{M.~S.~Vallabhajosyula, S.~S.~Budhya, and R.~Ramnath
\newblock ``Reference Implementation of Smart Scheduler: A CI-Aware, AI-Driven Scheduling Framework for HPC Workloads''.
\newblock {\em Practice and Experience in Advanced Research Computing 2024: Human Powered Computing}, pp. 1--4, 2024.}


\chapter*{Conference Posters \& Presentations}

\pubitem{M.~S.~Vallabhajosyula and R.~Ramnath
\newblock ``Towards Characterizing DNNs to Estimate Training Time using HARP (HPC Application Resource (runtime) Predictor)''.
\newblock {\em Practice and Experience in Advanced Research Computing 2023: Computing for the Common Good}, pp. 483--485, 2023.}

\pubitem{M.~S.~Vallabhajosyula, S.~S.~Budhya, A.~Jain, M.~Baig, and R.~Ramnath
\newblock ``Orchestrating a DNN Training Job Using an iScheduler Framework: A Use Case''.
\newblock {\em Practice and Experience in Advanced Research Computing 2024: Human Powered Computing}, pp. 1--3, 2024.}




% \pubitem{B.~Simpson
% \newblock ``Milking a Cow''.
% \newblock {\em Journal of Dairy Science}, 00(2):277--287, Feb. 1900.}



% \instructpubs
%
% \pubitem{B.~Simpson, ed.,
% \newblock ``Lab notes for Cow Science 101'', 1909.}

\end{publist}



\begin{fieldsstudy}

% The \majorfield* uses the unit specified in the \unit command used
% earlier in your document. If you want to use a different unit, use the
% second form shown here
% \majorfield*
\majorfield{Artificial Intelligence for High Performance Computing Optimization}

%%
%% Note:  If there were only one field of study, the following list 
%%        would best be done using the following command:
%%
%%  \onestudy{Only Topic}{Only Professor}
%%

% \begin{studieslist}
% \studyitem{Topic 1}{Prof.\ Big Dude}
% \studyitem{Topic 2}{Prof.\ Other Dude}
% \studyitem{Topic 3}{Prof.\ Another Dude}
% \end{studieslist}

\end{fieldsstudy}

\end{vita}

